%AttackName: M-Saddle Point Attack-1

%Input:
\( f \) be the target model
\( x \) be the original input 
\( y \) be the true label
\( \epsilon \) be the maximum allowable perturbation
\( \delta \) generate adversarial perturbation such that the model is misled by maximizing the loss, while the defender minimizes it.

%Output:
The output is a perturbation \( \delta \) that forms a saddle point of the loss function \( L(f(x + \delta), y) \), satisfying both attack and defense objectives.

%Formula:
\min_{\theta}\max_{\|\delta\|\le\epsilon} L\bigl(f_{\theta}(x+\delta),y\bigr)

%Explanation:
The Saddle Point Adversarial Attack exploits the interaction between the adversary and the defender. The adversary maximizes the loss by perturbing the input \( x \) within the allowable range \( \|\delta\| \leq \epsilon \), while the defender minimizes the loss by adjusting the model parameters \( \theta \). This creates a dynamic process that seeks a saddle point in the loss landscape, demonstrating the vulnerabilities of the model when exposed to adversarial perturbations.